% titel: Blatt 0, A1: Proportionalität

% Teil I: Werkzeuge, kein Lösungsweg.

% Proportionalität: die Grundidee des "Skalierens". Warum der Quotient
% Y/X konstant ist, was die Konstante bedeutet und wie man am Datensatz
% erkennt, ob echte Proportionalität vorliegt.

Die direkt proportionale Abhängigkeit ist der einfachste, aber auch der
grundlegendste funktionale Zusammenhang: $Y$ entsteht aus $X$, indem man
$X$ immer mit \emph{derselben} Zahl multipliziert. Anschaulich bedeutet
das nicht mehr und nicht weniger als \emph{Skalieren}.

\begin{werkzeugkasten}[Direkte Proportionalität: Grundidee]
$Y=k\,X$ heisst: Verdoppelt man $X$, verdoppelt sich auch $Y$; verdreifacht
man $X$, verdreifacht sich $Y$. Es gibt \emph{eine feste} Zahl $k$, die
alle Wertepaare zugleich verbindet. Zwei Kennzeichen, an denen man sie
sofort erkennt:
\begin{itemize}
  \item \textbf{Ursprungsgerade:} Der Graph von $Y$ gegen $X$ ist eine
        Gerade durch den Nullpunkt. Zu $X=0$ gehört stets $Y=0$ --- wer das
        nicht sieht, übersieht den Kern der Proportionalität.
  \item \textbf{Konstanter Quotient:} Der Quotient $k=Y/X$ ist für
        \emph{jedes} Paar derselbe. Diese beinahe triviale Beobachtung ist
        das wichtigste Werkzeug der Aufgabe.
\end{itemize}
\end{werkzeugkasten}

\begin{werkzeugkasten}[Warum spricht man von ``proportional'', nicht nur ``linear''?]
Eine lineare Funktion $Y=kX+d$ mit $d\neq0$ ist \emph{nicht} proportional,
denn zu $X=0$ gehört $Y=d\neq0$. Proportionalität ist der Spezialfall
$d=0$: Es gibt keine ``Grundlast'', die Grössen hängen vollständig und nur
über den Faktor $k$ zusammen. Hat also ein realer Vorgang eine feste
Zusatzkomponente (Eintrittspreis pro Besuch plus fixe Monatsgebühr), ist
er linear, aber nie proportional.
\end{werkzeugkasten}

\begin{werkzeugkasten}[Prüfen auf Proportionalität]
\begin{enumerate}
  \item Für jedes Wertepaar den Quotienten $k=\tfrac{Y}{X}$ bilden.
  \item Sind \emph{alle} Quotienten gleich, liegt direkte Proportionalität
        vor und $k$ ist dieser gemeinsame Quotient.
  \item Ist ein Quotient verschieden (oder $X=0$ vorgegeben), ist die
        direkte Proportionalität mit fester Konstante ausgeschlossen.
\end{enumerate}
\textbf{Warum nicht nur ein Paar, sondern alle?} Ein einzelnes Paar liefert
immer irgendeinen Quotienten. Die Aussage ``$Y$ verhält sich proportional
zu $X$'' ist aber eine Aussage über \emph{alle} Daten; schon ein einziges
abweichendes Paar widerlegt sie vollständig.
\end{werkzeugkasten}

\begin{achtung}[Darauf achten]
\begin{itemize}
  \item Die Konstante als gekürzten Bruch angeben, z.\,B. $k=\tfrac32$
        statt $1{,}5$; prüfen, dass jeder Quotient denselben Wert ergibt.
  \item Nicht nur ein Paar betrachten --- proportional heisst für
        \emph{alle} Paare derselbe Faktor.
  \item Division durch $0$ ausschliessen: Ist $X=0$ im Spiel, ist
        $k=Y/X$ nicht definiert.
  \item Nicht ``linear`` mit ``proportional`` verwechseln: Ein
        Achsenabschnitt ungleich $0$ zerstört die Proportionalität.
\end{itemize}
\end{achtung}
